\chapter{Estructuras de Lewis, resonancia y modelo RPECV}\label{ch:b1:lewis-resonance-vsepr}

La molécula de ozono, \ce{O3}, son tres átomos de oxígeno en una cadena
angular. Si se dibuja con las reglas del libro 1, un enlace sale doble y
el otro sencillo: un enlace debería ser corto, y el otro, largo. Sin
embargo, la espectroscopia de microondas encuentra los dos enlaces
\ce{O-O} del ozono exactamente iguales, de \qty{127.8}{pm}, entre el
enlace doble de \ce{O2} (\qty{120.8}{pm}) y el enlace sencillo del
peróxido de hidrógeno (\qty{147.5}{pm}). Ningún dibujo por sí solo es
correcto; dos dibujos juntos sí lo son. Este capítulo precisa el modelo
de Lewis --- \omterm{def:b1:lewis-resonance-vsepr:formal-charge}{cargas formales}, octetos incompletos o superados, resonancia
entre varias estructuras --- y lo usa después para predecir la forma de
las moléculas y los \omterm{def:b1:lewis-resonance-vsepr:dipole}{momentos dipolares} que de ella se derivan.
% ledger: geo:O3.rOO, re:O2, geo:H2O2.rOO

\begin{recall}
El libro 1 (año 10) describió el \omterm{def:b1:lewis-resonance-vsepr:lewis}{enlace covalente} como un par de
electrones compartido, dibujó \omterm{def:b1:lewis-resonance-vsepr:lewis}{estructuras de Lewis} con \omterm{def:b1:lewis-resonance-vsepr:lewis}{pares enlazantes} y
no enlazantes, y predijo cuatro formas (lineal, angular, trigonal y
tetraédrica). La \omterm{def:b1:periodicity:electronegativity}{electronegatividad} y sus diferencias se hicieron
cuantitativas en el \cref{ch:b1:periodicity}; los \omterm{def:b1:quantum-numbers:configuration}{electrones de valencia}
de un átomo se leen en su configuración
(\cref{def:b1:quantum-numbers:configuration}).
\end{recall}

\section{Estructuras de Lewis y cargas formales}

\begin{definition}[Enlace covalente, estructura de Lewis]\label{def:b1:lewis-resonance-vsepr:lewis}
Un \emph{enlace covalente}\index{enlace covalente} es un par de electrones
compartido por dos átomos, el \emph{par enlazante}\index{par enlazante};
dos o tres pares forman un enlace doble o triple. Un par de \omterm{def:b1:quantum-numbers:configuration}{electrones de
valencia} que pertenece a un solo átomo es un \emph{par no
enlazante}\index{par no enlazante}. Una \emph{estructura de
Lewis}\index{estructura de Lewis} muestra todos los \omterm{def:b1:quantum-numbers:configuration}{electrones de valencia}
de una molécula o de un ion, como enlaces (trazos entre átomos) y pares no
enlazantes (trazos o pares de puntos sobre un átomo). Según la \emph{regla
del octeto}\index{regla del octeto}, los átomos del \omterm{def:b1:periodicity:table}{periodo} 2 (C, N, O, F)
están rodeados en una estructura estable por cuatro pares, ocho
electrones; según la \emph{regla del dúo}\index{regla del duo@regla del dúo}, el hidrógeno, por un
par.
\end{definition}

\begin{definition}[Carga formal]\label{def:b1:lewis-resonance-vsepr:formal-charge}
La \emph{carga formal}\index{carga formal} de un átomo en una \omterm{def:b1:lewis-resonance-vsepr:lewis}{estructura
de Lewis} es
\[
  q_F = N_v - N_{\text{ne}} - \tfrac12 N_{\text{enl}},
\]
donde $N_v$ es el número de \omterm{def:b1:quantum-numbers:configuration}{electrones de valencia} del átomo libre,
$N_{\text{ne}}$ el número de electrones de sus \omterm{def:b1:lewis-resonance-vsepr:lewis}{pares no enlazantes} y
$N_{\text{enl}}$ el número de electrones de los enlaces que forma. Cuando
no es nula, se escribe $\oplus$ o $\ominus$ junto al átomo.
\end{definition}

\begin{proposition}[Las cargas formales suman la carga]\label{prop:b1:lewis-resonance-vsepr:formal-sum}
La suma de las \omterm{def:b1:lewis-resonance-vsepr:formal-charge}{cargas formales} de los átomos de una \omterm{def:b1:lewis-resonance-vsepr:lewis}{estructura de Lewis}
es igual a la carga total de la molécula o del ion.
\end{proposition}

\begin{proof}
Al sumar las \omterm{def:b1:lewis-resonance-vsepr:formal-charge}{cargas formales} $q_F$ de todos los átomos de la estructura se
obtiene $\sum N_v - (\text{electrones
no enlazantes} + \text{electrones de enlace})$, porque los
electrones de cada enlace se cuentan a medias en cada uno de sus dos
átomos. El paréntesis es el número total de \omterm{def:b1:quantum-numbers:configuration}{electrones de valencia}
dibujados en la estructura, que es igual a $\sum N_v$ menos la carga $z$
de la especie. Por tanto, $\sum q_F = z$.
\end{proof}

\begin{method}[Dibujar una estructura de Lewis]\label{met:b1:lewis-resonance-vsepr:lewis}
Para dibujar la \omterm{def:b1:lewis-resonance-vsepr:lewis}{estructura de Lewis} de una molécula o de un ion:
\begin{enumerate}
  \item se cuentan los \omterm{def:b1:quantum-numbers:configuration}{electrones de valencia} $N = \sum N_v - z$ y los
        pares $N/2$;
  \item se dibuja el esqueleto: el átomo menos electronegativo (salvo H)
        suele ser el central; el hidrógeno y el flúor son siempre
        terminales;
  \item se completan los octetos de los átomos terminales con \omterm{def:b1:lewis-resonance-vsepr:lewis}{pares no
        enlazantes} y se colocan los pares restantes en el átomo central;
  \item si al átomo central le falta el octeto, se convierten \omterm{def:b1:lewis-resonance-vsepr:lewis}{pares no
        enlazantes} de sus vecinos en enlaces múltiples;
  \item se calculan las \omterm{def:b1:lewis-resonance-vsepr:formal-charge}{cargas formales}; entre las estructuras posibles se
        prefiere la que tiene menos \omterm{def:b1:lewis-resonance-vsepr:formal-charge}{cargas formales} y las cargas negativas
        sobre los átomos más electronegativos.
\end{enumerate}
\end{method}

\begin{example}[El ácido nítrico]\label{ex:b1:lewis-resonance-vsepr:hno3}
\ce{HNO3}: $N = 1 + 5 + 3 \times 6 = 24$ electrones, 12 pares. Esqueleto:
nitrógeno central, unido a tres oxígenos, uno de los cuales lleva el
hidrógeno. Al completar los octetos con enlaces sencillos, al nitrógeno le
quedan seis electrones; un \omterm{def:b1:lewis-resonance-vsepr:lewis}{par no enlazante} de un oxígeno se convierte en
un enlace \ce{N=O}. \omterm{def:b1:lewis-resonance-vsepr:formal-charge}{Cargas formales}: nitrógeno $5 - 0 - 4 = +1$; el
oxígeno con enlace sencillo y sin hidrógeno, $6 - 6 - 1 = -1$; los demás,
0. El total es 0, como exige la
\cref{prop:b1:lewis-resonance-vsepr:formal-sum}.
\end{example}

\begin{omfigure}
\setchemfig{atom sep=2.4em}
\begin{tikzpicture}[font=\small]
  \node at (0,0) {\large\chemfig{H-\charge{90=\|,270=\|}{O}-\charge{90:2pt=$\scriptstyle\oplus$}{N}(=[:45]\charge{0=\|,90=\|}{O})-[:-45]\charge{0=\|,240=\|,300=\|,45:3pt=$\scriptstyle\ominus$}{O}}};
  \node[below] at (0,-1.5) {ácido nítrico, \ce{HNO3}};
  \node at (5.2,0) {\large\chemfig{\charge{180=\|}{N}~\charge{0=\|}{N}}};
  \node[below] at (5.2,-1.5) {dinitrógeno, \ce{N2}};
  \node at (9.6,0) {\large\chemfig{\charge{135=\|,225=\|,90:2pt=$\scriptstyle\ominus$}{N}=\charge{90:2pt=$\scriptstyle\oplus$}{N}=\charge{45=\|,315=\|}{O}}};
  \node[below] at (9.6,-1.5) {monóxido de dinitrógeno, \ce{N2O}};
\end{tikzpicture}

{\small \omterm{def:b1:lewis-resonance-vsepr:lewis}{Estructuras de Lewis} con \omterm{def:b1:lewis-resonance-vsepr:lewis}{pares no enlazantes} y \omterm{def:b1:lewis-resonance-vsepr:formal-charge}{cargas formales}.}
\end{omfigure}

\section{Más allá del octeto}

\begin{definition}[Moléculas deficientes en electrones e hipervalentes]\label{def:b1:lewis-resonance-vsepr:hypervalent}
Una molécula en la que un átomo tiene menos de ocho \omterm{def:b1:quantum-numbers:configuration}{electrones de
valencia} es \emph{deficiente en electrones}\index{deficiente en electrones}: en \ce{BF3}, el
boro tiene seis. Una molécula en la que un átomo del \omterm{def:b1:periodicity:table}{periodo} 3 o
siguientes está rodeado por más de cuatro pares es
\emph{hipervalente}\index{hipervalente}: en \ce{PCl5} el fósforo forma
cinco enlaces, y en \ce{SF6} el azufre, seis. Las moléculas con un número
impar de electrones, como \ce{NO} y \ce{NO2}, no pueden cumplir la \omterm{def:b1:lewis-resonance-vsepr:lewis}{regla
del octeto} en todos sus átomos: un electrón queda desapareado.
\end{definition}

\begin{remark}[Por qué el periodo 3 y no el periodo 2]\label{rem:b1:lewis-resonance-vsepr:hypervalent}
El nitrógeno forma \ce{NF3} pero nunca \ce{NF5}, mientras que el fósforo
forma tanto \ce{PF3} como \ce{PF5}. Un átomo del \omterm{def:b1:periodicity:table}{periodo} 2 es demasiado
pequeño para unirse a más de cuatro vecinos; un átomo del \omterm{def:b1:periodicity:table}{periodo} 3 es lo
bastante grande para acoger cinco o seis. Se encuentran dibujos de
\ce{SO4^{2-}} o de \ce{H2SO4} con enlaces dobles en el azufre (12
electrones alrededor de S) o con enlaces sencillos y \omterm{def:b1:lewis-resonance-vsepr:formal-charge}{cargas formales} (un
octeto); el segundo refleja mejor la distribución de carga, y ambos dan
la geometría correcta.
\end{remark}

\begin{definition}[Ácido de Lewis, base de Lewis, enlace dativo]\label{def:b1:lewis-resonance-vsepr:lewis-acid}
Un \emph{ácido de Lewis}\index{acido de Lewis@ácido de Lewis} es una especie capaz de aceptar
un par de electrones en un hueco de su capa de valencia (\ce{BF3},
\ce{AlCl3}, \ce{H+}, cationes metálicos); una \emph{base de
Lewis}\index{base de Lewis} es una especie con un \omterm{def:b1:lewis-resonance-vsepr:lewis}{par no enlazante} que
puede compartir (\ce{NH3}, \ce{H2O}, \ce{F-}). Cuando la base cede el par
al ácido, el enlace formado es un \emph{enlace dativo}\index{enlace dativo}:
un \omterm{def:b1:lewis-resonance-vsepr:lewis}{enlace covalente} cuyos dos electrones proceden de uno solo de los
participantes.
\end{definition}

\begin{example}[El amoniaco y el trifluoruro de boro]\label{ex:b1:lewis-resonance-vsepr:adduct}
\ce{NH3 + BF3 -> H3NBF3}. En el aducto, el boro completa su octeto; el
nitrógeno, que ahora comparte su \omterm{def:b1:lewis-resonance-vsepr:lewis}{par no enlazante}, lleva la \omterm{def:b1:lewis-resonance-vsepr:formal-charge}{carga formal}
$+1$, y el boro, $-1$. Una vez formado, el enlace \ce{N-B} es un \omterm{def:b1:lewis-resonance-vsepr:lewis}{enlace
covalente} ordinario.
\end{example}

\begin{omfigure}
\setchemfig{atom sep=2.1em}
\begin{tikzpicture}[font=\small]
  \node at (0,0) {\large\chemfig{H-\charge{90=\|}{N}(-[:240]H)(-[:300]H)}};
  \node at (1.6,0) {$+$};
  \node at (3.1,0) {\large\chemfig{B(-[:90]F)(-[:210]F)(-[:330]F)}};
  \draw[-{Stealth}, thick] (4.3,0) -- (5.3,0);
  \node at (7.4,0) {\large\chemfig{H-\charge{45:1pt=$\scriptstyle\oplus$}{N}(-[:90]H)(-[:270]H)-\charge{45:1pt=$\scriptstyle\ominus$}{B}(-[:90]F)(-[:270]F)-F}};
\end{tikzpicture}

{\small Un \omterm{def:b1:lewis-resonance-vsepr:lewis-acid}{enlace dativo}: el \omterm{def:b1:lewis-resonance-vsepr:lewis}{par no enlazante} del amoniaco (una \omterm{def:b1:lewis-resonance-vsepr:lewis-acid}{base de
Lewis}) ocupa el hueco del trifluoruro de boro (un \omterm{def:b1:lewis-resonance-vsepr:lewis-acid}{ácido de Lewis}). No se
dibujan los \omterm{def:b1:lewis-resonance-vsepr:lewis}{pares no enlazantes} del flúor.}
\end{omfigure}

\section{Resonancia}

\begin{definition}[Formas resonantes, híbrido de resonancia]\label{def:b1:lewis-resonance-vsepr:resonance}
Cuando los electrones de una molécula pueden colocarse en varias
\omterm{def:b1:lewis-resonance-vsepr:lewis}{estructuras de Lewis} que solo difieren en la posición de los enlaces
$\pi$ y de los \omterm{def:b1:lewis-resonance-vsepr:lewis}{pares no enlazantes} --- sin que se muevan los núcleos ---,
cada una es una \emph{forma resonante}\index{forma resonante} (o forma
mesómera), unida a las demás por una flecha de doble punta
$\leftrightarrow$. La molécula real no es ninguna de ellas: es el
\emph{híbrido de resonancia}\index{hibrido de resonancia@híbrido de resonancia}, una única estructura
cuya distribución electrónica es una media ponderada de las suyas.
\end{definition}

\begin{remark}[Una flecha, no un equilibrio]\label{rem:b1:lewis-resonance-vsepr:arrow}
La flecha $\leftrightarrow$ no describe una reacción: el ozono no pasa de
una estructura a la otra. El híbrido es una sola molécula, descrita con
varios dibujos porque uno solo no basta. Nunca debe usarse
$\rightleftharpoons$ entre \omterm{def:b1:lewis-resonance-vsepr:resonance}{formas resonantes}.
\end{remark}

\begin{method}[Escribir y ponderar formas resonantes]\label{met:b1:lewis-resonance-vsepr:resonance}
Para hallar las \omterm{def:b1:lewis-resonance-vsepr:resonance}{formas resonantes} de una especie:
\begin{enumerate}
  \item se parte de una \omterm{def:b1:lewis-resonance-vsepr:lewis}{estructura de Lewis}; se busca un \omterm{def:b1:lewis-resonance-vsepr:lewis}{par no
        enlazante} o un enlace $\pi$ junto a un enlace $\pi$, un hueco o
        una carga positiva;
  \item se desplazan pares con flechas curvas (un \omterm{def:b1:lewis-resonance-vsepr:lewis}{par no enlazante} se
        convierte en un enlace $\pi$, un enlace $\pi$ en un \omterm{def:b1:lewis-resonance-vsepr:lewis}{par no
        enlazante}), sin mover nunca un núcleo ni superar el octeto de un
        átomo del \omterm{def:b1:periodicity:table}{periodo} 2;
  \item se recalculan las \omterm{def:b1:lewis-resonance-vsepr:formal-charge}{cargas formales};
  \item se ponderan las estructuras: las más importantes cumplen la \omterm{def:b1:lewis-resonance-vsepr:lewis}{regla
        del octeto}, tienen menos \omterm{def:b1:lewis-resonance-vsepr:formal-charge}{cargas formales} y sitúan las cargas
        negativas sobre los átomos más electronegativos; las estructuras
        equivalentes pesan lo mismo.
\end{enumerate}
\end{method}

\begin{example}[Ozono, nitrato, benceno]\label{ex:b1:lewis-resonance-vsepr:examples}
El ozono tiene dos formas equivalentes: cada enlace \ce{O-O} es doble en
una y sencillo en la otra, de modo que ambos tienen un orden de enlace
$1.5$ y la misma longitud, \qty{127.8}{pm}, entre \ce{O=O} y \ce{O-O}. El
ion nitrato \ce{NO3-} tiene tres formas equivalentes; cada enlace
\ce{N-O} tiene orden $4/3$. El benceno, \ce{C6H6}, tiene dos estructuras
de Kekulé; sus seis enlaces \ce{C-C} son iguales, de \qty{139.7}{pm},
entre el enlace sencillo del etano (\qty{153.6}{pm}) y el doble del eteno
(\qty{133.9}{pm}).
% ledger: geo:O3.rOO, geo:C6H6.rCC, geo:C2H6.rCC, geo:C2H4.rCC
\end{example}

\begin{omfigure}
\vspace*{10pt}
\large
\setchemfig{atom sep=2.4em}
\schemestart
\chemfig{\charge{135=\|,225=\|}{O}=[:-30]\charge{270=\|,90:2pt=$\scriptstyle\oplus$}{O}-[:30]\charge{90=\|,0=\|,270=\|,45:5pt=$\scriptstyle\ominus$}{O}}
\arrow{<->}
\chemfig{\charge{90=\|,180=\|,270=\|,135:5pt=$\scriptstyle\ominus$}{O}-[:-30]\charge{270=\|,90:2pt=$\scriptstyle\oplus$}{O}=[:30]\charge{45=\|,315=\|}{O}}
\schemestop

\medskip
\schemestart
\chemfig{*6(=-=-=-)}
\arrow{<->}
\chemfig{*6(-=-=-=)}
\schemestop
\hspace{3em}
\chemfig{\charge{270:2pt=$\scriptstyle\oplus$}{N}(=[:90]\charge{45=\|,135=\|}{O})(-[:210]\charge{150=\|,240=\|,300=\|,90:2pt=$\scriptstyle\ominus$}{O})-[:330]\charge{30=\|,300=\|,240=\|,90:2pt=$\scriptstyle\ominus$}{O}}
\normalsize
\par\vspace{14pt}
{\small Resonancia en el ozono (un \omterm{def:b1:lewis-resonance-vsepr:lewis}{par no enlazante} del oxígeno de la
derecha se convierte en un enlace $\pi$, mientras que el enlace $\pi$ de
la izquierda se convierte en un \omterm{def:b1:lewis-resonance-vsepr:lewis}{par no enlazante}) y en el benceno, y una
de las tres formas equivalentes del ion nitrato, cuya carga se reparte
entre los tres oxígenos.}
\end{omfigure}

\begin{definition}[Enlace $\sigma$, enlace $\pi$]\label{def:b1:lewis-resonance-vsepr:sigma-pi}
Un enlace sencillo es un \emph{enlace $\sigma$}\index{enlace sigma@enlace $\sigma$}:
su par de electrones está sobre el eje que une los dos núcleos, por el
solapamiento frontal de dos orbitales. El segundo y el tercer enlace de un
enlace múltiple son \emph{enlaces $\pi$}\index{enlace pi@enlace $\pi$}: sus
pares están a ambos lados del eje, por el solapamiento lateral de dos
orbitales p perpendiculares a él. Un enlace doble es un enlace $\sigma$ y
uno $\pi$; un enlace triple, un $\sigma$ y dos $\pi$.
\end{definition}

\begin{omfigure}
\begin{tikzpicture}[font=\small]
  % sigma: two lobes overlapping head-on
  \fill[omDef!35, opacity=0.8] (0,0) ellipse (0.9 and 0.42);
  \fill[omProp!35, opacity=0.8] (1.1,0) ellipse (0.9 and 0.42);
  \fill (0.1,0) circle (1.6pt); \fill (1.0,0) circle (1.6pt);
  \draw[dashed, black!50] (-1.2,0) -- (2.3,0);
  \node[below] at (0.55,-1.05) {$\sigma$: frontal, sobre el eje};
  % pi: two p orbitals side by side
  \begin{scope}[xshift=5.2cm]
    \foreach \x/\c in {0/omDef, 1.2/omProp} {
      \fill[\c!35, opacity=0.85] (\x,0.45) ellipse (0.32 and 0.45);
      \fill[\c!15, opacity=0.85] (\x,-0.45) ellipse (0.32 and 0.45);
      \draw[\c!70!black] (\x,0.45) ellipse (0.32 and 0.45) (\x,-0.45) ellipse (0.32 and 0.45);
      \fill (\x,0) circle (1.6pt);
    }
    \draw[dashed, black!50] (-0.7,0) -- (1.9,0);
    \draw[omThm, thick, densely dotted] (0.25,0.75) -- (0.95,0.75) (0.25,-0.75) -- (0.95,-0.75);
    \node[below] at (0.6,-1.05) {$\pi$: lateral, arriba y abajo};
  \end{scope}
\end{tikzpicture}

{\small Los dos tipos de solapamiento. Un enlace $\sigma$ permite la
rotación alrededor de su eje; un enlace $\pi$ no, y por eso un enlace
doble es rígido
(\cref{ch:b1:stereochemistry-in-depth}).}
\end{omfigure}

\section{El modelo RPECV}

\begin{definition}[Modelo RPECV, número estérico]\label{def:b1:lewis-resonance-vsepr:vsepr}
En el \emph{modelo RPECV}\index{modelo RPECV} (repulsión de los pares de
electrones de la capa de valencia), los pares que rodean un átomo central
A se repelen y adoptan la disposición que los mantiene más alejados entre
sí. Una molécula se escribe \ce{AX_nE_m}, donde $n$ es el número de
átomos unidos a A (un enlace múltiple cuenta una vez) y $m$ el número de
\omterm{def:b1:lewis-resonance-vsepr:lewis}{pares no enlazantes} de A. El \emph{número estérico}\index{numero esterico@número estérico} $n + m$
fija la disposición de los pares: 2, lineal; 3, trigonal plana; 4,
tetraédrica; 5, bipiramidal trigonal; 6, octaédrica. La forma de la
molécula es la disposición de sus átomos únicamente.
\end{definition}

\begin{proposition}[Geometrías RPECV]\label{prop:b1:lewis-resonance-vsepr:geometries}
El \omterm{def:b1:lewis-resonance-vsepr:vsepr}{modelo RPECV} predice las formas y los ángulos ideales de la tabla
siguiente; los ángulos medidos son próximos.
\begin{center}
\small
\begin{tabular}{llllc}
\hline
tipo & disposición de los pares & forma & ejemplo & ángulo medido \\
\hline
\ce{AX2} & lineal & lineal & \ce{CO2} & $180^\circ$ \\
\ce{AX3} & trigonal plana & trigonal plana & \ce{BF3} & $120^\circ$ \\
\ce{AX2E} & trigonal plana & angular & \ce{SO2} & $119.5^\circ$ \\
\ce{AX4} & tetraédrica & tetraédrica & \ce{CH4} & $109.5^\circ$ \\
\ce{AX3E} & tetraédrica & piramidal trigonal & \ce{NH3} & $106.7^\circ$ \\
\ce{AX2E2} & tetraédrica & angular & \ce{H2O} & $104.5^\circ$ \\
\ce{AX5} & bipiramidal trigonal & bipiramidal trigonal & \ce{PCl5} & $90^\circ$, $120^\circ$ \\
\ce{AX4E} & bipiramidal trigonal & de balancín & \ce{SF4} & $101.6^\circ$, $173.1^\circ$ \\
\ce{AX3E2} & bipiramidal trigonal & en forma de T & \ce{ClF3} & $87.5^\circ$ \\
\ce{AX2E3} & bipiramidal trigonal & lineal & \ce{XeF2} & $180^\circ$ \\
\ce{AX6} & octaédrica & octaédrica & \ce{SF6} & $90^\circ$ \\
\ce{AX5E} & octaédrica & piramidal cuadrada & \ce{BrF5} & --- \\
\ce{AX4E2} & octaédrica & plano-cuadrada & \ce{XeF4} & --- \\
\hline
\end{tabular}
\end{center}
% ledger: geo:CO2.aOCO, geo:BF3.aFBF, geo:SO2.aOSO, geo:CH4.aHCH,
% geo:NH3.aHNH, geo:H2O.aHOH, geo:SF6.aFSF, geo:SF4.aFSF2, geo:SF4.aFSF,
% geo:ClF3.aFClF, geo:XeF2.aFXeF
\end{proposition}
\admitted

\begin{proposition}[Los pares no enlazantes cierran los ángulos]\label{prop:b1:lewis-resonance-vsepr:lone-pair-angles}
Un \omterm{def:b1:lewis-resonance-vsepr:lewis}{par no enlazante}, sujeto por un solo núcleo, ocupa más espacio que un
\omterm{def:b1:lewis-resonance-vsepr:lewis}{par enlazante} y repele con más fuerza: las repulsiones disminuyen en el
orden \omterm{def:b1:lewis-resonance-vsepr:lewis}{par no enlazante}/\omterm{def:b1:lewis-resonance-vsepr:lewis}{par no enlazante} $>$ \omterm{def:b1:lewis-resonance-vsepr:lewis}{par no enlazante}/\omterm{def:b1:lewis-resonance-vsepr:lewis}{par
enlazante} $>$ \omterm{def:b1:lewis-resonance-vsepr:lewis}{par enlazante}/\omterm{def:b1:lewis-resonance-vsepr:lewis}{par enlazante}. Por eso los ángulos entre
enlaces se cierran a medida que se añaden \omterm{def:b1:lewis-resonance-vsepr:lewis}{pares no enlazantes}
($109.5^\circ$ en \ce{CH4}, $106.7^\circ$ en \ce{NH3}, $104.5^\circ$ en
\ce{H2O}), y en una bipirámide trigonal los \omterm{def:b1:lewis-resonance-vsepr:lewis}{pares no enlazantes} ocupan las
posiciones ecuatoriales, donde solo tienen dos vecinos a $90^\circ$.
% ledger: geo:CH4.aHCH, geo:NH3.aHNH, geo:H2O.aHOH
\end{proposition}
\admitted

\begin{omfigure}
\setchemfig{atom sep=2.6em}
\renewcommand{\arraystretch}{1.2}
\begin{tabular}{ccccc}
\chemfig{O=C=O} &
\chemfig{B(-[:90]F)(-[:210]F)(-[:330]F)} &
\chemfig{C(-[:90]H)(-[:200]H)(<[:280]H)(<:[:335]H)} &
\chemfig{P(-[:90]Cl)(-[:270]Cl)(-[:180]Cl)(<[:320]Cl)(<:[:30]Cl)} &
\chemfig{S(-[:90]F)(-[:270]F)(-[:180]F)(-[:0]F)(<[:225]F)(<:[:45]F)} \\
\small \ce{AX2}, \ce{CO2} & \small \ce{AX3}, \ce{BF3} & \small \ce{AX4}, \ce{CH4} &
\small \ce{AX5}, \ce{PCl5} & \small \ce{AX6}, \ce{SF6} \\[16pt]
\chemfig{\charge{90=\:}{S}(=[:210]O)(=[:330]O)} &
\chemfig{\charge{90=\:}{N}(-[:200]H)(<[:280]H)(<:[:335]H)} &
\chemfig{\charge{60=\:,120=\:}{O}(-[:215]H)(-[:325]H)} &
\chemfig{\charge{0=\:}{S}(-[:90]F)(-[:270]F)(<[:210]F)(<:[:150]F)} &
\chemfig{Xe(-[:0]F)(-[:90]F)(-[:180]F)(-[:270]F)} \\
\small \ce{AX2E}, \ce{SO2} & \small \ce{AX3E}, \ce{NH3} & \small \ce{AX2E2}, \ce{H2O} &
\small \ce{AX4E}, \ce{SF4} & \small \ce{AX4E2}, \ce{XeF4} \\
\end{tabular}

\medskip
{\small Formas predichas por el \omterm{def:b1:lewis-resonance-vsepr:vsepr}{modelo RPECV}. Una cuña apunta hacia el
lector, y una cuña rayada, en sentido contrario; los enlaces simples
están en el plano de la página. Los \omterm{def:b1:lewis-resonance-vsepr:lewis}{pares no enlazantes} del átomo central
se representan como pares de puntos; \ce{XeF4} se dibuja de frente, con
sus dos \omterm{def:b1:lewis-resonance-vsepr:lewis}{pares no enlazantes} apuntando hacia el lector y en sentido
contrario.}
\end{omfigure}

\begin{method}[Predecir una forma]\label{met:b1:lewis-resonance-vsepr:vsepr}
Para predecir la forma alrededor de un átomo central:
\begin{enumerate}
  \item se dibuja la \omterm{def:b1:lewis-resonance-vsepr:lewis}{estructura de Lewis} (\cref{met:b1:lewis-resonance-vsepr:lewis});
  \item se cuentan los átomos unidos al átomo central ($n$) y sus \omterm{def:b1:lewis-resonance-vsepr:lewis}{pares
        no enlazantes} ($m$); se escribe \ce{AX_nE_m};
  \item se lee la disposición a partir de $n + m$, y la forma, a partir de
        las posiciones de los $n$ átomos;
  \item se corrigen los ángulos ideales: los \omterm{def:b1:lewis-resonance-vsepr:lewis}{pares no enlazantes} y los
        enlaces múltiples ocupan más espacio que los enlaces sencillos.
\end{enumerate}
\end{method}

\begin{definition}[Hibridación]\label{def:b1:lewis-resonance-vsepr:hybridisation}
La \emph{hibridación}\index{hibridacion@hibridación} de un átomo es una
descripción de sus enlaces mediante orbitales que apuntan en las
direcciones predichas por el \omterm{def:b1:lewis-resonance-vsepr:vsepr}{modelo RPECV}: de un átomo de \omterm{def:b1:lewis-resonance-vsepr:vsepr}{número estérico}
4 se dice que es sp$^3$ (cuatro orbitales equivalentes a $109.5^\circ$);
de \omterm{def:b1:lewis-resonance-vsepr:vsepr}{número estérico} 3, sp$^2$ (tres orbitales a $120^\circ$ en un plano, y
un orbital p restante perpendicular a él); de \omterm{def:b1:lewis-resonance-vsepr:vsepr}{número estérico} 2, sp (dos
orbitales a $180^\circ$, y dos orbitales p restantes). Los orbitales p
restantes forman los enlaces $\pi$.
\end{definition}

\begin{example}[Los carbonos de la química orgánica]\label{ex:b1:lewis-resonance-vsepr:carbon}
En el etano cada carbono es sp$^3$, tetraédrico; en el eteno, sp$^2$, con
los seis átomos en un plano y ángulos próximos a $120^\circ$ (\ce{H-C-H}
medido: $117.6^\circ$); en el etino, sp, con los cuatro átomos en una
recta. Las longitudes de enlace se acortan de 153.6 a 133.9 y a
\qty{120.3}{pm} a medida que se añaden enlaces $\pi$.
% ledger: geo:C2H4.aHCH, geo:C2H6.rCC, geo:C2H4.rCC, geo:C2H2.rCC
\end{example}

\begin{remark}[Una descripción, no una causa]\label{rem:b1:lewis-resonance-vsepr:hybrid-limit}
La \omterm{def:b1:lewis-resonance-vsepr:hybridisation}{hibridación} describe una geometría ya conocida; no la explica, y
falla con las energías de los electrones, que explica la teoría de
orbitales moleculares, en el volumen del segundo año. Sigue siendo el
lenguaje cotidiano de la química orgánica: un ``carbono sp$^2$'' es un
carbono plano con un orbital p libre para un enlace $\pi$.
\end{remark}

\section{Moléculas polares: momentos dipolares}

\begin{definition}[Dipolo de enlace, momento dipolar]\label{def:b1:lewis-resonance-vsepr:dipole}
En un enlace entre átomos de \omterm{def:b1:periodicity:electronegativity}{electronegatividades} distintas, los
electrones se desplazan hacia el átomo más electronegativo, que lleva una
carga parcial $-\delta e$ mientras que su compañero lleva $+\delta e$
($0 < \delta < 1$). El \emph{dipolo de enlace}\index{dipolo de enlace} es
el vector $\vec\mu = \delta e\,
\vec d$, de norma $\delta e d$, que apunta de la carga parcial negativa a
la positiva, siendo $\vec d$ el vector que las une. El \emph{momento
dipolar}\index{momento dipolar} de una molécula es la suma vectorial de
sus dipolos de enlace (y de las contribuciones de sus \omterm{def:b1:lewis-resonance-vsepr:lewis}{pares no
enlazantes}); se expresa en culombios metro o en debyes, $\qty{1}{\debye} =
\qty{3.336e-30}{C.m}$. Una molécula con momento dipolar no nulo es una
\emph{molécula polar}\index{molecula polar@molécula polar}.
% ledger: const:c (1 D = 1e-21/c C.m)
\end{definition}

\begin{remark}[Convenio de sentido]\label{rem:b1:lewis-resonance-vsepr:convention}
La física dibuja el vector dipolar de la carga negativa a la positiva,
el convenio que se usa aquí. Los libros de química más antiguos dibujan
una flecha con la cola cruzada que apunta hacia el extremo negativo; solo
cambia el sentido.
\end{remark}

\begin{proposition}[Dipolo de una molécula simétrica]\label{prop:b1:lewis-resonance-vsepr:dipole-sum}
Una molécula \ce{AX_n} sin \omterm{def:b1:lewis-resonance-vsepr:lewis}{pares no enlazantes} en A (\ce{AX2} lineal,
\ce{AX3} trigonal, \ce{AX4} tetraédrica, \dots) tiene un \omterm{def:b1:lewis-resonance-vsepr:dipole}{momento dipolar}
nulo, sea cual sea la polaridad de sus enlaces. Una molécula \ce{AX2}
angular, de ángulo $\theta$, cuyos enlaces tienen el dipolo $\mu_b$, tiene
el \omterm{def:b1:lewis-resonance-vsepr:dipole}{momento dipolar}
\[
  \mu = 2\mu_b \cos\frac{\theta}{2} .
\]
\end{proposition}

\begin{proof}
En las formas simétricas, los vectores de enlace son de igual norma y sus
direcciones están dispuestas simétricamente alrededor de A, de modo que
suman cero (para \ce{AX2} lineal son opuestos; para \ce{AX3} a
$120^\circ$, tres vectores de la misma norma suman cero; para \ce{AX4} la
suma es invariante por las rotaciones del tetraedro, luego es nula). En
la molécula angular, cada vector de enlace forma el ángulo $\theta/2$ con
la bisectriz; las componentes perpendiculares a la bisectriz se anulan y
las paralelas a ella se suman: $2\mu_b\cos(\theta/2)$.
\end{proof}

\begin{omfigure}
\begin{tikzpicture}[font=\small, >={Stealth}]
  % water
  \node[aO] (o) at (0,0) {};
  \node[aH] (h1) at (-0.95,-0.73) {};
  \node[aH] (h2) at (0.95,-0.73) {};
  \begin{scope}[on background layer]
    \draw[bond] (o.center) -- (h1.center); \draw[bond] (o.center) -- (h2.center);
  \end{scope}
  \draw[->, omProp, thick] (-0.25,0.15) -- (-0.95,-0.38);
  \draw[->, omProp, thick] (0.25,0.15) -- (0.95,-0.38);
  \draw[->, omThm, very thick] (0,-0.25) -- (0,-1.45);
  \node[omThm, right] at (0.05,-1.25) {$\vec\mu$};
  \node at (0,0.55) {$\delta^-$};
  \node at (-1.25,-1.05) {$\delta^+$}; \node at (1.25,-1.05) {$\delta^+$};
  \node[below] at (0,-1.7) {\ce{H2O}: $\mu = \qty{1.86}{\debye}$};
  % carbon dioxide
  \begin{scope}[xshift=5cm]
    \node[aO] (a) at (-1.1,0) {}; \node[aC] (c) at (0,0) {}; \node[aO] (b) at (1.1,0) {};
    \begin{scope}[on background layer]
      \foreach \d in {-1.6,1.6} {
        \draw[bond, line width=1.1pt] ([yshift=\d pt]a.center) -- ([yshift=\d pt]c.center);
        \draw[bond, line width=1.1pt] ([yshift=\d pt]c.center) -- ([yshift=\d pt]b.center);}
    \end{scope}
    \draw[->, omProp, thick] (-0.75,0.45) -- (-0.2,0.45);
    \draw[->, omProp, thick] (0.75,0.45) -- (0.2,0.45);
    \node[below] at (0,-1.7) {\ce{CO2}: los dipolos de enlace se anulan, $\mu = 0$};
  \end{scope}
\end{tikzpicture}

{\small \omterm{def:b1:lewis-resonance-vsepr:dipole}{Dipolos de enlace} (en naranja, de la carga parcial negativa a la
positiva) y su suma (en rojo). En el agua, angular, se suman; en el
dióxido de carbono, lineal, se anulan.}
% ledger: dip:H2O
\end{omfigure}

\begin{example}[Moléculas polares y apolares]\label{ex:b1:lewis-resonance-vsepr:dipoles}
\ce{CO2}, \ce{BF3}, \ce{CH4} y \ce{CCl4} tienen \omterm{def:b1:lewis-resonance-vsepr:dipole}{momento dipolar} nulo.
\ce{H2O} (\qty{1.86}{\debye}), \ce{NH3} (\qty{1.48}{\debye}) y
\ce{SO2} (\qty{1.63}{\debye}) son polares. A lo largo de \ce{CH3Cl},
\ce{CH2Cl2} y \ce{CHCl3}, el momento disminuye, 1.87, 1.62 y
\qty{1.04}{\debye}, hasta anularse en \ce{CCl4}: los \omterm{def:b1:lewis-resonance-vsepr:dipole}{dipolos de enlace}
\ce{C-Cl} se compensan cada vez más. En los haluros de hidrógeno sigue la
diferencia de \omterm{def:b1:periodicity:electronegativity}{electronegatividad}: HF 1.83, HCl 1.09, HBr 0.83 y HI
\qty{0.45}{\debye}.
% ledger: dip:H2O, dip:NH3, dip:SO2, dip:CH3Cl, dip:CH2Cl2, dip:CHCl3,
% dip:CCl4, dip:HF, dip:HCl, dip:HBr, dip:HI
\end{example}

\section{Ejercicios}

\begin{exercise}[$\star$]\label{exo:b1:lewis-resonance-vsepr:1}
Dibújense las \omterm{def:b1:lewis-resonance-vsepr:lewis}{estructuras de Lewis} de \ce{H2O}, \ce{NH3}, \ce{CO2},
\ce{HCN}, \ce{CH2O} (metanal) y \ce{NH4+}, e indíquense las \omterm{def:b1:lewis-resonance-vsepr:formal-charge}{cargas
formales}.
\end{exercise}

\begin{exercise}[$\star$]\label{exo:b1:lewis-resonance-vsepr:2}
Dense el tipo \ce{AX_nE_m}, la forma y el ángulo aproximado de \ce{H2S},
\ce{PCl3}, \ce{BF3}, \ce{SiH4} y \ce{CO3^{2-}} (el átomo central es el
primero de cada fórmula).
\end{exercise}

\begin{exercise}[$\star$]\label{exo:b1:lewis-resonance-vsepr:3}
¿Cuáles de estas moléculas son polares: \ce{CCl4}, \ce{CH2Cl2}, \ce{BF3},
\ce{NH3}, \ce{CO2}, \ce{SO2}? Justifíquese a partir de las formas.
\end{exercise}

\begin{exercise}[$\star$]\label{exo:b1:lewis-resonance-vsepr:4}
Identifíquense el \omterm{def:b1:lewis-resonance-vsepr:lewis-acid}{ácido de Lewis} y la \omterm{def:b1:lewis-resonance-vsepr:lewis-acid}{base de Lewis} en
\ce{H+ + H2O -> H3O+} y en \ce{Cu^{2+} + 4NH3 -> [Cu(NH3)4]^{2+}}. ¿Qué
enlaces son dativos?
\end{exercise}

\begin{exercise}[$\star\star$]\label{exo:b1:lewis-resonance-vsepr:5}
Escríbanse dos \omterm{def:b1:lewis-resonance-vsepr:resonance}{formas resonantes} del ion etanoato \ce{CH3COO-}. ¿Qué
predicen para las longitudes de los dos enlaces \ce{C-O}? Compárese con
el ácido metanoico \ce{HCOOH}, cuyos dos enlaces \ce{C-O} miden 120.2 y
\qty{134.3}{pm}.
% ledger: geo:H2CO2.rCO, geo:H2CO2.rCO2
\end{exercise}

\begin{exercise}[$\star\star$]\label{exo:b1:lewis-resonance-vsepr:6}
Dibújese la \omterm{def:b1:lewis-resonance-vsepr:lewis}{estructura de Lewis} del ion sulfato \ce{SO4^{2-}} con cuatro
enlaces sencillos \ce{S-O}, calcúlense las \omterm{def:b1:lewis-resonance-vsepr:formal-charge}{cargas formales} y predígase la
forma. ¿Cuántas \omterm{def:b1:lewis-resonance-vsepr:resonance}{formas resonantes} con dos enlaces \ce{S=O} y sin carga
en el azufre pueden escribirse?
\end{exercise}

\begin{exercise}[$\star\star$]\label{exo:b1:lewis-resonance-vsepr:7}
Predíganse con el \omterm{def:b1:lewis-resonance-vsepr:vsepr}{modelo RPECV} las formas de \ce{SF4}, \ce{ClF3},
\ce{XeF2} y \ce{XeF4}. Explíquese por qué los \omterm{def:b1:lewis-resonance-vsepr:lewis}{pares no enlazantes} de una
bipirámide trigonal se sitúan en el plano ecuatorial.
\end{exercise}

\begin{exercise}[$\star\star$]\label{exo:b1:lewis-resonance-vsepr:8}
Los ángulos \ce{H-X-H} valen $104.5^\circ$ en \ce{H2O} y $92.1^\circ$ en
\ce{H2S}; $106.7^\circ$ en \ce{NH3} y $93.3^\circ$ en \ce{PH3}.
Propóngase una explicación basada en el tamaño y la \omterm{def:b1:periodicity:electronegativity}{electronegatividad}
del átomo central.
% ledger: geo:H2O.aHOH, geo:H2S.aHSH, geo:NH3.aHNH, geo:PH3.aHPH
\end{exercise}

\begin{exercise}[$\star\star$]\label{exo:b1:lewis-resonance-vsepr:9}
Dese la \omterm{def:b1:lewis-resonance-vsepr:hybridisation}{hibridación} de cada carbono y del nitrógeno en \ce{CH3-C#N}
(etanonitrilo) y en \ce{CH2=CH-CH3}, y el número de enlaces $\sigma$ y
$\pi$ de cada molécula.
\end{exercise}

\begin{exercise}[$\star\star\star$]\label{exo:b1:lewis-resonance-vsepr:10}
El \omterm{def:b1:lewis-resonance-vsepr:dipole}{momento dipolar} de \ce{NF3} es de solo \qty{0.24}{\debye}, frente a
\qty{1.48}{\debye} para \ce{NH3}, aunque el enlace \ce{N-F} es más polar
que el enlace \ce{N-H}. Explíquese la diferencia a partir de los \omterm{def:b1:lewis-resonance-vsepr:dipole}{dipolos
de enlace} y del \omterm{def:b1:lewis-resonance-vsepr:lewis}{par no enlazante} del nitrógeno.
% ledger: dip:NF3, dip:NH3
\end{exercise}

\begin{exercise}[$\star\star\star$]\label{exo:b1:lewis-resonance-vsepr:11}
La molécula de \ce{SO2} tiene un ángulo de $119.5^\circ$ y un \omterm{def:b1:lewis-resonance-vsepr:dipole}{momento
dipolar} de \qty{1.63}{\debye}. Calcúlese el dipolo de un enlace \ce{S-O};
con la longitud de enlace, \qty{143.2}{pm}, dedúzcase la carga parcial
$\delta$ de cada oxígeno.
% ledger: geo:SO2.aOSO, geo:SO2.rSO, dip:SO2
\end{exercise}

\begin{exercise}[$\star\star\star$]\label{exo:b1:lewis-resonance-vsepr:12}
Dibújense las \omterm{def:b1:lewis-resonance-vsepr:lewis}{estructuras de Lewis} más importantes del monóxido de
dinitrógeno \ce{N2O} (lineal, \ce{N-N-O}) y calcúlense las \omterm{def:b1:lewis-resonance-vsepr:formal-charge}{cargas
formales}. Las longitudes de enlace medidas son \ce{N-N} \qty{112.8}{pm} y
\ce{N-O} \qty{118.4}{pm}; en \ce{N2} el enlace triple mide
\qty{109.8}{pm}. ¿Qué estructura pesa más?
% ledger: geo:N2O.rNN, geo:N2O.rNO, re:N2
\end{exercise}

\section{Problema: las moléculas de un rayo}

\begin{problem}[{Problema de fin de semana --- las especies del nitrógeno
y del oxígeno que se forman cuando un rayo calienta el aire: estructuras
de Lewis, resonancia, formas y cargas parciales del agua}]
\label{pb:b1:lewis-resonance-vsepr:1}
Un rayo calienta el aire a miles de grados: \ce{N2} y \ce{O2} dan
\ce{NO}, después \ce{NO2}, \ce{N2O}, ozono y, por último, ácido nítrico en
la lluvia. Datos: $\qty{1}{\debye} = \qty{3.336e-30}{C.m}$, $e =
\qty{1.602e-19}{C}$. Los valores medidos se dan donde hacen falta.

\medskip
\noindent\textbf{Parte I --- \omterm{def:b1:lewis-resonance-vsepr:lewis}{Estructuras de Lewis}.}
\begin{enumerate}
  \item Cuéntense los \omterm{def:b1:quantum-numbers:configuration}{electrones de valencia} de \ce{N2}, \ce{NO},
        \ce{NO2}, \ce{N2O}, \ce{O3} y \ce{HNO3}.
  \item Dibújese la \omterm{def:b1:lewis-resonance-vsepr:lewis}{estructura de Lewis} de \ce{N2}.
  \item Dibújese una \omterm{def:b1:lewis-resonance-vsepr:lewis}{estructura de Lewis} de \ce{NO}. ¿Por qué no puede
        cumplirse la \omterm{def:b1:lewis-resonance-vsepr:lewis}{regla del octeto} en los dos átomos?
  \item Dibújense dos \omterm{def:b1:lewis-resonance-vsepr:lewis}{estructuras de Lewis} de \ce{N2O}, con \ce{N=N=O} y
        con \ce{N#N-O}, e indíquense las \omterm{def:b1:lewis-resonance-vsepr:formal-charge}{cargas formales}.
  \item Dibújese una \omterm{def:b1:lewis-resonance-vsepr:lewis}{estructura de Lewis} del ozono e indíquense las
        \omterm{def:b1:lewis-resonance-vsepr:formal-charge}{cargas formales}.
  \item Dibújese una \omterm{def:b1:lewis-resonance-vsepr:lewis}{estructura de Lewis} de \ce{HNO3} e indíquense las
        \omterm{def:b1:lewis-resonance-vsepr:formal-charge}{cargas formales}.
  \item ¿Cuáles de las seis especies tienen un \omterm{def:b1:quantum-numbers:unpaired}{electrón desapareado}?
\end{enumerate}

\medskip
\noindent\textbf{Parte II --- Resonancia y longitudes de enlace.}
\begin{enumerate}[resume]
  \item Escríbanse las dos \omterm{def:b1:lewis-resonance-vsepr:resonance}{formas resonantes} del ozono. ¿Qué orden de
        enlace predicen?
  \item El enlace \ce{O-O} mide \qty{127.8}{pm} en el ozono,
        \qty{120.8}{pm} en \ce{O2} y \qty{147.5}{pm} en \ce{H2O2}. ¿Se
        confirma la predicción?
  \item Escríbanse las tres \omterm{def:b1:lewis-resonance-vsepr:resonance}{formas resonantes} de \ce{NO3-} e indíquese el
        orden de cada enlace \ce{N-O}.
  \item En \ce{HNO3} los tres enlaces \ce{N-O} miden 140.6, 121.1 y
        \qty{119.9}{pm}. Asígnense y explíquese el resultado con la
        resonancia.
  \item Explíquese por qué el ángulo \ce{O-N-O} de \ce{NO2} ($134^\circ$)
        es mayor que $120^\circ$.
  \item ¿Cuál de las dos estructuras de la pregunta 4 sitúa la \omterm{def:b1:lewis-resonance-vsepr:formal-charge}{carga
        formal} negativa sobre el átomo más electronegativo?
\end{enumerate}

\medskip
\noindent\textbf{Parte III --- Formas y polaridad.}
\begin{enumerate}[resume]
  \item Dese el tipo \ce{AX_nE_m} del átomo central de \ce{O3},
        \ce{N2O} y \ce{NO3-} y del nitrógeno de \ce{HNO3}, así como sus
        formas.
  \item El ángulo del ozono es de $116.8^\circ$. Explíquese por qué es
        menor que $120^\circ$.
  \item El ozono tiene un \omterm{def:b1:lewis-resonance-vsepr:dipole}{momento dipolar} de \qty{0.53}{\debye}; \ce{N2O},
        de \qty{0.16}{\debye}; \ce{N2}, ninguno. Explíquese.
  \item ¿Por qué \ce{CO2} es apolar mientras que \ce{SO2} es polar?
  \item Predígase el ángulo de \ce{NH3} respecto de $109.5^\circ$ y
        compárese después con el valor medido, $106.7^\circ$.
  \item La misma pregunta para el agua ($104.5^\circ$).
\end{enumerate}
% ledger: geo:O3.rOO, re:O2, geo:H2O2.rOO, geo:HNO3.rNO, geo:HNO3.rNO2,
% geo:HNO3.rNO3, geo:NO2.aONO, geo:O3.aOOO, dip:O3, dip:N2O

\medskip
\noindent\textbf{Parte IV --- Las cargas parciales del agua.}
El \omterm{def:b1:lewis-resonance-vsepr:dipole}{momento dipolar} del agua es de \qty{1.857}{\debye}; su ángulo, de
$104.48^\circ$, y la longitud de su enlace \ce{O-H}, de \qty{95.8}{pm}.
% ledger: dip:H2O, geo:H2O.aHOH, geo:H2O.rOH
\begin{enumerate}[resume]
  \item Exprésese el \omterm{def:b1:lewis-resonance-vsepr:dipole}{momento dipolar} del agua en función del \omterm{def:b1:lewis-resonance-vsepr:dipole}{dipolo de
        enlace} $\mu_{OH}$ y del ángulo.
  \item Calcúlese $\mu_{OH}$ en debyes.
  \item Conviértase en culombios metro.
  \item Si el \omterm{def:b1:lewis-resonance-vsepr:dipole}{dipolo de enlace} estuviera formado por dos cargas enteras
        $\pm e$ situadas en los dos núcleos, ¿cuánto valdría?
  \item Dedúzcase el carácter iónico fraccionario $\delta$ del enlace
        \ce{O-H}, la razón entre el \omterm{def:b1:lewis-resonance-vsepr:dipole}{dipolo de enlace} medido y el de cargas
        enteras.
\end{enumerate}
\end{problem}
